\hypertarget{native-291-remarks}{}\pdfbookmark[1]{Bemerkungen: Einzelbemerkungen}{bemerkungen-remarks}
Zu Seite 45, Zeile 1 findet sich eine Notiz von \textsc{Gauss} in seinem Handexemplar,
wonach dort $\omega'$ statt $\omega$ stehen soll; indessen ist $\omega$ richtig, falls in
Gleichung XI., Art. 22, nur das zweite Glied mit Hülfe hyperbolischer Logarithmen berechnet
wird, das erste ($e\operatorname{tang} F$) aber durch briggische; $\omega'$ muss stehen,  wenn \textls[150]{beide}
Glieder mit hyperbolischen Logarithmen berechnet werden; dann muss aber $\lambda$ im
Nenner fortfallen.

Zu den numerischen Beispielen ist im Speciellen folgendes zu bemerken: Auf den Planeten
\textls[150]{Juno} beziehen sich ausser Art. 150---155 auch die Zahlen in den Art. 10, 13, 14, 51,
63---65, 69, 77, 87\,I, 97\,I und die Berechnung des Sonnenorts Art. 73; auf den Planeten
\textls[150]{Pallas} ausser Art. 156---157 auch Art. 56 und 58 und die Berechnung des Sonnenorts
Art. 70; auf den Planeten \textls[150]{Ceres} ausser Art. 159 auch Art. 87\,II, 97\,II. Ausserdem stehen
im Zusammenhang die Beispiele in den Art. 38, 43, 97\,III für eine Ellipse mit grosser
Excentricität, und Art. 23, 24, 26, 46, 105 für eine Hyperbel.

Zu Art. 150---155: Die Beobachtungen der Juno teilt \textsc{Maskelyne} in einem Brief an
\textsc{Gauss} vom 8. Mai 1805 mit. Die Sonnenörter, sowie die Nutation und die Schiefe der
Ekliptik hat \textsc{Gauss} aus den \textsc{Zach}schen Sonnentafeln entnommen; auf eine
genaue Controlle derselben (Seite 195, Zeile 8---10) wurde verzichtet, da ihre strenge
Richtigkeit für die Durchführung der Beispiele unwesentlich ist und sie mit Rücksicht auf
die weitere Rechnung auch nicht hätten geändert werden können; doch ist Seite 195, Zeile 8
im ersten Sonnenort $53''{,}72$ in $53''{,}71$ verwandelt worden, um das Folgende hiemit zur
Übereinstimmung zu bringen. Die Werthe der Constanten, welche \textsc{Gauss} zur Reduction
der Beobachtungen benutzt hat, sind
\[
\begin{aligned}
\text{Präcession}&=50''{,}11\\
\text{Aberration}&=20''{,}00\\
\text{Längendifferenz Paris-Greenwich}&=9^{\mathrm m}20^{\mathrm s}{,}8.
\end{aligned}
\]

Die Werthe für die Elemente (Seite 204) ergaben sich bei der Controllrechnung theilweise
etwas abweichend; die Genauigkeitsgrenze, innerhalb welcher sie überhaupt zu bestimmen
sind, hängt von der Schärfe ab, mit der die Grösse $x = \sin\tfrac12 g^2$ bekannt ist. In
der ersten Hypothese wurde bei der \SourceLeaf{292}Controllrechnung gefunden
\[
\tfrac12 g = 1^0 34' 10''{,}046 \qquad x = 0{,}000\,7501\,446
\]
und hieraus ergab sich:
\[
\begin{aligned}
\text{Epoche} &= 42^0 12' 39''{,}89 \\
\text{mittl. Bewegung} &= 823''{,}4987 \\
\text{Perihel} &= 52^0 41' 13''{,}24 \\
\varphi &= 14\ 24\ 28{,}58 \\
\text{Logarithm. der halben grossen Axe} &= 0{,}422\,8958 \\
\text{Aufsteigender Knoten} &= 171^0 5' 48''{,}84 \\
\text{Neigung} &= 13^0 2' 37''{,}50.
\end{aligned}
\]

Die von \textsc{Gauss} angegebenen Werthe scheinen auf der Annahme
\[
\tfrac12 g = 1^0 34' 10''{,}056 \qquad x = 0{,}000\,7501\,474
\]
zu beruhen. Die Differenzen liegen also unterhalb der zu erreichenden Genauigkeitsgrenze.
In der zweiten Hypothese wurden sie noch geringer gefunden, in der dritten ergaben sich
genau die \textsc{Gauss}schen Werthe, abgesehen von den im Fehlerverzeichniss vermerkten
Stellen.

Zu Art. 156---157. Die Declination der dritten Beobachtung der Pallas ist in
\textsc{Zach}s Monatlicher Correspondenz a.\,a.\,O. zu $-28^0 38' 8''{,}3$ angegeben; die
Sonnenörter Seite 205, Zeile 5---7 von unten sind nicht controllirt. Die Nutation hat
\textsc{Gauss} vermuthlich nach seinen eigenen Tafeln gerechnet, und zwar nicht nach den in
der Monatlichen Correspondenz (s. Band VI, Seite 123) abgedruckten, sondern nach
handschriftlich in ein Handbuch (B\,b) eingetragenen, die als Vorarbeit von den gedruckten
sich nur durch andere Werthe der Constanten unterscheiden; die Nutation fand sich richtig
bis auf den dritten Werth ($13''{,}68$ statt $12''{,}20$), die Aberration mit dem Werth der
Constanten $20''{,}00$ gerechnet, bis auf den dritten und vierten ($21''{,}51$ und $1''{,}63$ statt
$23''{,}68$ und $1''{,}31$). Die Präcession hat \textsc{Gauss} mit den Werthen $m = 45''{,}93$,
$n = 20''{,}02$ gerechnet; hier differirt der dritte Ort ($2''{,}55$ statt $2''{,}25$). Diese
Differenzen haben beim Abdruck nur theilweise berücksichtigt werden können, vgl. das
Fehlerverzeichniss.

Die Elemente Seite 211 fanden sich z.\,Th. auch etwas abweichend, aber innerhalb der
Genauigkeitsgrenze der ganzen Rechnung, z.\,B. Perihel vom Knoten: $323^0 44' 56''{,}64$.

Zu Art. 159. Die Beobachtungszeit des ersten Orts der Ceres ist im Briefe \textsc{Olbers}
an \textsc{Gauss} vom 8. September 1805 zu $13^{\mathrm h} 8^{\mathrm m} 47^{\mathrm s}$
angegeben; der Reduction sind die Längendifferenzen Bremen-Paris $= 25^{\mathrm m}
51^{\mathrm s}$, Göttingen-Paris $= 30^{\mathrm m} 21^{\mathrm s}$, Lilienthal-Paris $=
26^{\mathrm m} 14^{\mathrm s}$ zum Grunde gelegt, welche in \textsc{Zach}, Tabulae motuum
Solis, Supplementum, Gotha 1804, mitgetheilt sind. Seite 213 Zeile 8---11 sind die Werthe
für Schiefe der Ekliptik, Präcession und Nutation und Zeile 9 u. 8 von unten die
Sonnenörter keiner Controlle unterworfen worden. Seite 214 haben sich in den
Zahlenangaben für Hypothese II grössere Unrichtigkeiten herausgestellt, die nur zum Theil
verbessert werden konnten. Die Werthe von $\omega+\sigma$, $\log Qc\sin\omega$ und $z$
sind verbessert worden; der darauf folgende Werth von $r'$ ergibt sich dann auch so, wie
\textsc{Gauss} ihn angibt, ist also wahrscheinlich mit dem eingesetzten richtigen Werthe
von $z$ gerechnet; dagegen ist von hier ab die weitere Rechnung mit dem unrichtigen
Werthe der alten Ausgabe ($z = 7^0 2' 32''{,}870$) richtig weitergeführt Seite 215 ist die
Rechnung bis zum $\log r' = 0{,}413\,2817$ richtig; von hier ab aber scheint
\textsc{Gauss} mit den unrichtigen Werthen $\log\dfrac{n'r'}{n} = 0{,}662\,0412$ (statt
$0{,}662\,0380$) und $\log\dfrac{n'r'}{n''} = 0{,}636\,4081$ (statt $0{,}636\,4049$)
gerechnet zu haben, denen die weitere Rechnung richtig entspricht. Auch die Werthe für die
wahren Anomalien, den Winkel $\varphi$ und die halbe grosse Axe Seite 215, Zeile 11, 10,
5, 4 v. unten ergaben sich bei der Nachrechnung etwas abweichend, sowohl mit Anwendung des
in der alten Ausgabe Art. 97 gegebenen
\SourceLeaf{293}
Werthes $x = 0{,}065\,29078$, als auch des verbesserten $x = 0{,}065\,29076$; den von
\textsc{Gauss} angegebenen Werthen würde $g = 29^0 36' 31''{,}975$ ($x = 0{,}065\,290798$)
entsprechen.

Zu Art. 171. In der ersten Beobachtung der \textls[150]{Vesta} gibt \textsc{Olbers} im Briefe an
\textsc{Gauss} vom 5. April 1807 die Rectascension zu $183^0 52' 37''$ an. Seite 231 sind
die Sonnenörter Zeile 12---15, sowie die Ableitung der Zahlen ebenda Zeile 9---6 v. unten
aus den Beobachtungen nicht controllirt. Seite 235 scheinen die Werthe der Elemente mit
Hülfe von $g = 23^0 34' 22''{,}34$ ($x = 0{,}041\,7240$) gerechnet zu sein. Die Zahlen Seite
235, Zeile 7 bis 2 v. unten wurden ebenfalls nicht controllirt.

Alle handschriftlichen Bemerkungen, welche \textsc{Gauss} in sein Handexemplar der
Theoria motus eingetragen hat, sind als Fussnoten eingefügt und zwar Seite 20, 30, 52,
53, 56, 65, 66, 80, 117---118, 120, 144, 145, 183, 226, 242, 249, 251, 270, wobei in der
Notiz Seite 118, letzte Zeile, $\dfrac{17420}{729}$ in $\dfrac{15420}{729} =
\dfrac{5120}{243}$ und Seite 185, Zeile 7 v. unten $-3 \pm \BareRoot{(9-16\operatorname{cotang} A^2)}$ in
$-3\operatorname{cotang} A \pm \BareRoot{(9\operatorname{cotang} A^2-16)}$ verbessert worden ist. Die Entwickelung des
Ausdrucks in der Fussnote S. 270 gibt $y = 1 + \dfrac{10}{9}h + \dfrac{110}{81}hh +
\cdots$; es ist also vermuthlich (vgl. S. 118) ein Zeichenfehler begangen worden, und
sollte etwa heissen: $y = \left(1+\dfrac{32}{9}h\right)^{\frac{5}{16}}$. Zu Seite 144
vgl. die Bemerkungen über \textsc{Ivory}s Methode zur Cometenbahnbestimmung im Briefe
\textsc{Gauss} an \textsc{Olbers} vom 31. December 1814.

Ferner mag Folgendes bemerkt werden:

Zu Art. 32\,IV: Der Maximalfehler in der wahren Anomalie tritt für den Werth von $v$ ein,
der sich aus der Gleichung $\cos v + e\cos 2v = 0$ oder $\cos v = \dfrac1e\left\{
-\tfrac14 + \BareRoot{\left(\tfrac1{16}+\tfrac12 ee\right)}\right\}$ bestimmt; für alle auf
Seite 43 angegebenen Werthe von $e$ ist $v$ ungefähr gleich $60^0$.

Zu Art. 32\,VII: In der Hyperbel tritt der Maximalfehler für $F$ ein, wenn $\cos F +
\tfrac13\sin 2F - e\cos 2F = 0$. Die der Tabelle auf Seite 44 entsprechenden Werthe von
$F$ sind:

\begin{center}
\begin{tabular}{c|c}
$e$ & $F$ \\
\hline
$1{,}3$ & $14^0 46'$ \\
$1{,}2$ & $11\ 16$ \\
$1{,}1$ & $6\ 41$ \\
$1{,}05$ & $3\ 44$ \\
$1{,}01$ & $0\ 50$ \\
$1{,}001$ & $0\ 5$ \\
\end{tabular}
\end{center}

Zu Art. 36. Der Maximalfehler in der Ellipse tritt ein für $1-2ee+e\cos E = 0$, in der
Hyperbel für $ee(1+6\sin F)+e(\cos F-3\operatorname{tang} F) = 2+3\sin F$; die zu den Seite 48
gegebenen Tabellen gehörigen Werthe von $e$ sind also die folgenden:

\begin{center}
\begin{minipage}[t]{.42\linewidth}\centering
Ellipse:\par\smallskip
\begin{tabular}{c|c}
$E$ & $e$ \\
\hline
$10^0$ & $0{,}99494$ \\
$20$ & $0{,}98003$ \\
$30$ & $0{,}95602$ \\
$40$ & $0{,}92409$ \\
$50$ & $0{,}88583$ \\
$60$ & $0{,}83534$ \\
\end{tabular}
\end{minipage}\hfill
\begin{minipage}[t]{.42\linewidth}\centering
Hyperbel:\par\smallskip
\begin{tabular}{c|c}
$F$ & $e$ \\
\hline
$10^0$ & $1{,}00508$ \\
$20$ & $1{,}02094$ \\
$30$ & $1{,}04993$ \\
$40$ & $1{,}09754$ \\
$50$ & $1{,}17671$ \\
$60$ & $1{,}32007$ \\
$70$ & $1{,}63435$ \\
\end{tabular}
\end{minipage}
\end{center}
\SourceLeaf{294}
\textsc{Gauss}’ Vorarbeiten zur Theoria motus finden sich hauptsächlich in den mit B\,b, B\,c, B\,d bezeichneten Handbüchern, die Tafeln im Handbuch B\,b S. 30, 31, 35. Zur Entstehung des Manuskriptes und für die Vorbereitungen zur Drucklegung vergleiche man den Briefwechsel mit \textsc{Olbers} (W. \textsc{Olbers}, sein Leben und seine Werke, herausgegeben von C. \textsc{Schilling}, Band II).

\begin{flushright}
\textsc{Brendel}.
\end{flushright}

\begin{center}
\rule{0.3\linewidth}{0.4pt}
\end{center}
